\input{../tex-inputs/latex-leseansicht-vorspann}

\section[Gustav Schwarzkopf an Arthur Schnitzler, 21. 8. 1899]{L04293 Gustav Schwarzkopf an Arthur Schnitzler, 21. 8. 1899}
\nopagebreak\mylabel{L04293v}
\rehead{ }\normalsize\beginnumbering\briefempfaengerindex{Schnitzler, Arthur@\textsc{Schnitzler, Arthur}!zzzSchwarzkopf, Gustav@\emph{von Gustav Schwarzkopf}!1899-08-202@{20. 8. 1899}|(be}
\toendnotes[C]{\smallbreak\pagebreak[2]}
\correspDesc{Versand  durch Gustav Schwarzkopf am 20. 8. 1899 in Wien
\newline{}Übermittlung  am 21. 8. 1899 in Wien
\newline{}Zustellung  am 21. 8. 1899 in Bad Ischl
\newline{}Erhalt  durch Arthur Schnitzler am 21. 8. 1899 in Bad Ischl}\toendnotes[C]{\smallbreak}
\Standort{CUL, Schnitzler, B 96a.}
\physDesc{Postkarte, 571 Zeichen
\newline{}Handschrift: schwarze Tinte, deutsche Kurrent
\newline{}Versand: 1) Stempel: »\nobreak{}\oindex{I., Innere Stadt@\textbf{I., Innere Stadt}, \emph{Verwaltungsgebiet}|pwk}Wien 1/1 1, 21. 8. 99, 8–9V\nobreak{}«.   2) Stempel: »\nobreak{}\oindex{Bad Ischl@\textbf{Bad Ischl}|pwk}Ischl, 21. 8. 99, 11–12\nobreak{}«. }\toendnotes[C]{\smallbreak}\pstart{}{\pb}Herrn \textsc{D\textsuperscript{r} Arthur Schnitzler}\pend{}\pstart{}\textsc{Ischl\oindex{Bad Ischl@\textbf{Bad Ischl}|pw}}\pend{}\pstart{}\textsc{Pension Petter\oindex{Hotel und Pension Rudolfshöhe (Leopold Petter)@\textbf{Hotel und Pension Rudolfshöhe (Leopold Petter)}, \emph{Hotel}|pw}}\pend{}{\bigskip}\vspace{1em}
\pstart
           \raggedleft{}{\pb}20. 8. 99\pend
           \vspace{0.5em}
\pstart
           Lieber Arthur, den »\label{K_L04293-1v}\edtext{Ruck}{\lemma{\textnormal{\emph{Ruck}}}\Cendnote{\textnormal{XXXX Auszeichnungsfehler: Dokument L04203 nicht gefunden.}}}\label{K_L04293-1}« habe ich empfangen und \uline{geſpürt}; we{\geminationn} er mich doch
               nicht vom Fleck bringt,{ }ſo kö{\geminationn}en Sie erſehen, wie groß
               meine Unluſt iſt, die übrigens nebenbei auch von nüchternen praktiſchen Erwägungen
               geſtützt wird. Wir werden ja darüber{ }ſprechen. Wie lange bleiben Sie{ }ſelbſt de{\geminationn} noch in \textsc{Ischl\oindex{Bad Ischl@\textbf{Bad Ischl}|pw}}? Wol bis Ende des Monats? Und da{\geminationn}? Bleibt es bei
               der Stadtour nach Thüringen\oindex{Thüringen@\textbf{Thüringen}, \emph{Land}|pw}? Hier iſt es übrigens
               seit zwei Tagen{ }ſo kühles Wetter eher geeignet von Reiſegedanken »\label{K_L04293-2v}\edtext{zu entwöhnen als dazu anzureizen\pwindex{Schiller, Friedrich von 10.\,11.\,1759 Marbach am Neckar – 9.\,5.\,1805 Weimar@\textsc{Schiller, Friedrich von} (10.\,11.\,1759 Marbach am Neckar – 9.\,5.\,1805 Weimar), \emph{Schriftsteller, Historiker, Philosoph}!Kabale und Liebe@\strich\emph{Kabale und Liebe}|pwv}}{\lemma{\textnormal{\emph{zu … anzureizen}}}\Cendnote{\textnormal{Friedrich Schiller\pwindex{Schiller, Friedrich von 10.\,11.\,1759 Marbach am Neckar – 9.\,5.\,1805 Weimar@\textsc{Schiller, Friedrich von} (10.\,11.\,1759 Marbach am Neckar – 9.\,5.\,1805 Weimar), \emph{Schriftsteller, Historiker, Philosoph}|pwk}, \emph{Kabale und Liebe}\pwindex{Schiller, Friedrich von 10.\,11.\,1759 Marbach am Neckar – 9.\,5.\,1805 Weimar@\textsc{Schiller, Friedrich von} (10.\,11.\,1759 Marbach am Neckar – 9.\,5.\,1805 Weimar), \emph{Schriftsteller, Historiker, Philosoph}!Kabale und Liebe@\strich\emph{Kabale und Liebe}|pwk}, 4. Akt, 3. Szene: »Und
                     mit diesem ihr Herz zu teilen? – Ungeheuer! Unverantwortlich! – Einem Kerl,
                     mehr gemacht, von Sünden zu entwöhnen als dazu anzureizen.«
               }}}\label{K_L04293-2}«.\pend
           
\pstart
           Herzlichſt{\\[\baselineskip]}Ihr{\\[\baselineskip]}\spacefill\mbox{Gustav}\pend
           \leftskip=0em{}\selectlanguage{ngerman}\endnumbering\briefempfaengerindex{Schnitzler, Arthur@\textsc{Schnitzler, Arthur}!zzzSchwarzkopf, Gustav@\emph{von Gustav Schwarzkopf}!1899-08-202@{20. 8. 1899}|)be}\mylabel{L04293h}
\begin{anhang}
\end{anhang}\newcommand{\dateiname}{L04293}\newcommand{\titel}{Gustav Schwarzkopf an Arthur Schnitzler, 21. 8. 1899}\newcommand{\editorInnen}{Herausgegeben von Jahnke, SelmaMüller, Martin Anton}\input{../tex-inputs/latex-leseansicht-abspann}
